\section{\centering Solutions}

\subsection{Solutions to T/F Questions}
	\begin{table}[!h]
		\centering
		\begin{tabular}{c@{\hspace{2.5em}}cp{0.8\textwidth}}
			true & false & \\
			\specialrule{\lightrulewidth}{0.25em}{1.25em}
		         \bubble & \dbubble & Every line in \(\mathbb{R}^n\) a subspace.\\[0.75em]
			 \dbubble & \bubble & If \(A \vec{x} = A \vec{y} = \vec{b}\), for some \(A \in \mathbb{R}^{m \times n}\) and 
						       \(\vec{x}, \vec{y} \in \mathbb{R}^{n}\), then \(2\vec{x} - 2\vec{y} \in \nullspace A\).\\[0.75em]
			\bubble & \dbubble & The zero vector is part of any basis of a subspace of \(\mathbb{R}^n\). \\[0.75em]
			\dbubble & \bubble & Given a \(m \times n\) matrix \(A\) such that for all vectors \(\vec{x}, \vec{y} \in \mathbb{R}^n\), 
						      \(A\vec{x} = A\vec{y}\), then \(\nullspace A = \mathbb{R}^n\).\\[1.75em]
			\bubble & \dbubble & Given any \(m \times n\) matrix \(A\) and it's echelon form \(U\), \(\colspace A = \colspace U\).\\[0.75em]
			\dbubble & \bubble & Given any \(m \times n\) matrix \(A\) and it's echelon form \(U\), \(\nullspace A = \nullspace U\).\\[0.75em]
			\bubble & \dbubble & If \(A\) is \(n \times n\) and \(\rank A = n\), then \(\nullspace A\) is empty.\\[0.75em]
			\dbubble & \bubble & If \(\vec{v}\) lives in a 2-dimensional subspace of \(\mathbb{R}^3\), then it's coordinate under a basis for the subspace 
						      has 2 entries.\\[1.75em]
		\bottomrule
		\end{tabular}
	\end{table}

\subsection{Solutions to Possible/Impossible Questions}
	\begin{table}[!h]
		\centering
		\begin{tabular}{c@{\hspace{2.5em}}cp{0.7\textwidth}}
			possible & impossible & \\
			\specialrule{\lightrulewidth}{0.25em}{1.25em}
			\dbubble & \bubble & A matrix whose columns do not form a basis for its column space. \\[0.75em]
			\bubble & \dbubble & A 3-dimensional subpsace \(H\) of \(\mathbb{R}^4\), where the number of vectors in a basis for \(H\) is 2.\\[1.75em]
			\dbubble & \bubble & A vector living in 2-dimensional subspace of \(\mathbb{R}^6\), but there are only 2 entries for its coordinate vector.\\[1.75em]
			\bubble & \dbubble & An invertible \(n \times n\) with the dimension of the column space being less than \(n\). \\[1.75em]
			\dbubble & \bubble & The nullspace of a \(3 \times 2\) matrix that has 2 pivots has a dimension of 0. \\[0.75em]
			\bottomrule 
		\end{tabular}
	\end{table}

\subsection{Solutions to Short Answer Questions}
	\begin{enumerate}
		\item \([\vec{x}]_\mathcal{B} = \begin{bmatrix}3 \\ 2\end{bmatrix}\). \(\vec{x} = \begin{bmatrix}5 \\ 2\end{bmatrix}\).
		\item \[\begin{bmatrix}1 & 6 & 4 \\ 1 & 3 & 0\end{bmatrix} \sim \begin{bmatrix}1 & 6 & 4 \\ 0 & -3 & -4\end{bmatrix}.\] 
			Since there's a pivot in the first two columns, the first two columns of \(\begin{bmatrix}1 & 6 & 4 \\ 1 & 3 & 0\end{bmatrix}\), 
			\(\left\{\begin{bmatrix}1 \\ 0\end{bmatrix}, \begin{bmatrix}6 \\ 3\end{bmatrix}\right\}\) form \emph{a} basis for 
			\(\colspace \begin{bmatrix}1 & 6 & 4 \\ 1 & 3 & 0\end{bmatrix}\).
		
			Continuing off, 
			\[\begin{bmatrix}1 & 6 & 4 \\ 0 & -3 & -4\end{bmatrix} \sim \begin{bmatrix}1 & 0 & -4 \\ 0 & -3 & -4\end{bmatrix}
			  \sim \begin{bmatrix}1 & 0 & -4 \\ 0 & 1 & 4/3\end{bmatrix}\]
			This corresponds to the system
			 \begin{align*}
				      x_1 &        - 4x_3       = 0 \implies x_1 = 4x_3 \\
				      x_2 &+\frac{4}{3}x_3 = 0 \implies x_2 = -\frac{4}{3}x_3
			\end{align*}
			So 
			\[\vec{x} = \begin{bmatrix}x_1 \\ x_2 \\ x_3\end{bmatrix} = \begin{bmatrix}4x_3 \\ -4/3x_3 \\ x_3\end{bmatrix} 
				      = x_3\begin{bmatrix}4 \\ -4/3 \\ 1\end{bmatrix}.\]
			Notice what this is saying. It's saying that every \(\vec{x} \in \nullspace\biggl(\begin{bmatrix}1 & 6 & 4 \\ 1 & 3 & 0\end{bmatrix}\biggr)\) is a
			scalar multiple of \(\begin{bmatrix}4 \\ -4/3 \\ 1\end{bmatrix}\). That means that \(\begin{bmatrix}4 \\ -4/3 \\ 1\end{bmatrix}\) spans the null
			space of the matrix, and since it forms a linearly independent set, it is \emph{a} basis for the nullspace of this matrix.
	
	\end{enumerate}
